\subsection{Componentes incrementales de las cinco innovaciones}
La rama 1.9 conserva las cinco innovaciones de SD13 y sus entregables de diseño, construcción, pruebas específicas y evaluación. Cada código terminal de origen mantiene su correspondencia; cuando describe un período superior a una quincena se generan hijos por resultado y período. El esfuerzo se suma desde los hijos, sin agregar otra estimación al padre. Los componentes de IN4 e IN5 reciben paquetes terminales propios sin alterar los nombres de sus innovaciones.

La descomposición cuantifica 84 paquetes y 2.680 HH: IN1, 624; IN2, 712; IN3, 568; IN4, 352; IN5, 424. De ellas, 2.352 corresponden al alcance inicial con mes fijo y 328 al piloto estacional de IN2, ahora situado en 30--37 con invernada en 31. Los complementos de seguimiento y botadura se suman en el catálogo integrado. Estas horas no incluyen la revisión previa de cada embarcación ni el servicio recurrente descritos a continuación. Las horas de aceptación final de innovaciones del mes 21 son 120 y pertenecen al cierre de implementación; su soporte posterior se registra en operación. Los pilotos de IN2 desde $M_{\mathrm{inv}}$ pertenecen a la fase de operación y no se anticipan al precio de implementación.

Los paquetes específicos incluyen la evidencia de privacidad, accesibilidad, seguridad, continuidad y uso de IA que corresponde a la innovación. La construcción obligatoria de esos servicios comunes permanece en 1.2, 1.3 y 1.7; las pruebas integradas y la certificación contractual permanecen en 1.10. Una evidencia específica podrá alimentar el expediente general, sin imputar dos veces la misma ejecución \parencite[caps.~19, 20 y 26]{bases-transversales}.

\subsection{Revisión previa a botaduras y operación incremental}
\paragraph{1.9.2.5.4.k: revisión de pendientes previa a botadura}
Este paquete conserva el origen \texttt{IN2.5.4}. Se seleccionan cinco lotes $k=1,\ldots,5$ de cuatro embarcaciones: dos en 33, dos en 34 y uno en 35, para veinte embarcaciones. La distribución se ajustará al calendario confirmado de botaduras, conservando tres semanas de anticipación y sin duplicar el lote. La fecha de cada revisión deberá ser tres semanas anterior a su botadura prevista, conforme a SD13. El entregable es una lista por embarcación con plan aceptado, trabajos, pendientes, responsable de decisión y evidencia disponible. Excluye ejecutar mantenimiento físico, autorizar una navegación o asumir una certificación técnica ajena a Synaptix.

José Lara Arce responde por el paquete; Funcional revisa el plan y Calidad su evidencia. Esfuerzo: Funcional 4, Operación 8 y Calidad 4 HH; total 16. Se requieren planes, registros del piloto, calendario de botaduras confirmado, contratistas y validadores del cliente. Se recibirá con una revisión por cada embarcación del lote, diferencias identificadas y comunicación de decisiones pendientes; un pendiente sin resolver no equivale a autorización de botadura. Se estima una revisión documental consolidada por embarcación con antecedentes accesibles, no una inspección física ni una investigación de incidentes.

Para $N_{o,q}$ embarcaciones que requieren revisión, se define $B_{o,q}=\lceil N_{o,q}/4\rceil$ si $N_{o,q}>0$, y cero cuando el calendario confirmado no contiene revisiones. El esfuerzo es $16\sum B_{o,q}$ HH. La falta de calendario o antecedentes no justifica declarar cero: constituye una habilitación pendiente. Cambios de fecha obligan a actualizar la revisión y dimensionar el trabajo adicional. Predecesoras: planes aceptados e informes de seguimiento disponibles; sucesoras: resultados de botaduras y evaluación IN2.5.6. Hito: evaluación de la campaña, sin nuevo porcentaje de pago de implementación. Referencias: SD13, IN2.5.4, y ventanas del caso 07.

\paragraph{1.12.2.i.m.q: servicio quincenal de IN1, IN2 e IN3}
Se define un informe por innovación $i=1,2,3$, mes $m$ y quincena $q$. Los orígenes son \texttt{IN1.7.m.q}, \texttt{IN2.7} y \texttt{IN3.7}. José Lara Arce responde por el resultado. IN1 registra solicitudes, avisos fallidos e incidencias: Operación 4, Funcional 2 y Calidad 2 HH. IN2 utiliza un expediente mensual de 16 HH: Operación 8, Funcional 4 y Calidad 4, código 1.12.2.2.m. Conserva el mismo esfuerzo equivalente a dos informes de 8 HH y sus evidencias quincenales, sin sumar ambos formatos. IN3 registra procesamiento, consumo, formatos y fallos: Operación 4, Datos 2 y Calidad 2 HH. Los paquetes quincenales IN1/IN3 tienen 8 HH.

Se recibirá con versión del servicio, universo observado, incidencias, consumo o cambios aplicables, acciones y responsables trazados. Requiere telemetría, registros protegidos y servicio habilitado. Excluye guardias continuas, NOC/SOC, soporte base, nuevos desarrollos y los informes de piloto: estos últimos se imputan sólo a 1.9.2.5. IN1 e IN3 generan 72 paquetes cada una entre 21 y 56: 576 HH por innovación. IN2 genera $(57-M_{\mathrm{inv}})$ expedientes mensuales desde su habilitación estacional hasta 56: $16(57-M_{\mathrm{inv}})$ HH. Referencias: SD13, operación de cada innovación; contrato de operación y SLA. El volumen real de incidencias se contrastará para dimensionar soporte adicional.

\paragraph{1.12.2.4.m: expediente mensual de conciliación de resultados}
Conserva el trabajo recurrente de \texttt{IN4.4}, separado de la habilitación \texttt{1.9.4.4.1}. Ulysses Barreto responde por el expediente: Funcional 8, Calidad 4 y Proyecto 4 HH, total 16. El expediente conserva muestra, intervalos aceptados, cálculo versionado, condiciones de calidad, factor resultante, observaciones y decisión contractual. La aprobación del incentivo corresponde al cliente y no se sustituye por la preparación del informe.

Se recibirá con cálculo reproducible y vínculo a todas sus entradas; sin evidencia suficiente, muestra admisible o mejora atribuible el factor será cero. Se requieren registros validados de IN1, regla aceptada y responsables de preparación y aprobación distintos. No incluye de nuevo la captura del piloto, cobro automático ni un ahorro ya atribuido a IN1. Se genera un paquete mensual en 21--56: 36 paquetes, 576 HH. La espera de aprobación se programa aparte. Referencias: SD13, IN4.4, y E-25 para la separación de pagos de operación.

\paragraph{1.12.2.5.m: informe mensual de seguimiento ambiental}
Conserva el trabajo recurrente de \texttt{IN5.4}, separado de la habilitación \texttt{1.9.5.4.1}. Nicolás Torfan responde por el resultado: Datos 8, Calidad 4 y Operación 4 HH, total 16. Entregará cobertura de lecturas, permanencias conciliadas, consentimiento, acciones y comparación ambiental válida. Se recibirá con denominadores, períodos y exclusiones trazados; datos incompletos o incomparables impedirán declarar una reducción.

Requiere medición individual validada, historial, participantes autorizados y telemetría. Excluye compra y obras de medidores a cargo del cliente, captura base de consumo y mantenimiento físico. Se genera un paquete mensual en 21--56: 36 paquetes, 576 HH. El informe previo a la temporada 2029 se obtiene de esta serie y sus evidencias; no añade otra vez su esfuerzo de seguimiento. Referencias: SD13, IN5.4; SD2, EXC-06/EXC-10 y SUP-04; caso 07, capítulos 10 y 13.

Los cuatro servicios con calendario fijo suman 2.304 HH de operación; IN2 añade $16(57-M_{\mathrm{inv}})$ y la revisión de botaduras añade $16\sum B_{o,q}$. Los expedientes mensuales IN4/IN5 se preparan y controlan en la segunda quincena a partir de registros continuos del servicio; la espera de cierre, validación o aprobación del mes se enlaza aparte. Su frecuencia mensual no convierte el paquete de 16 HH en una tarea activa durante todo el mes. Son cargas incrementales por resultado, sin equivalencia a cobertura presencial ni atención continua. Se agregarán por persona a los paquetes generales de operación y a la estabilización de datos, conservando una imputación por tarea.
